\section{Exact output}\label{sec:exact}

A certified gap below $2^{-q}$ is not an exact solution. For rational
quadratics, a global minimizer and the optimal value have bounded rational
height, and this makes a sufficiently small certified gap convertible into an
exact optimizer that is checked against the original data. This section
proves the height bound in a form that serves two purposes at once, gives an
acceptance rule that needs neither uniqueness nor growth, and shows that an
exact optimizer can be recovered from any feasible point that is close to the
\emph{set} of minimizers. Under quadratic growth at a unique minimizer this
yields exact output within the bound of Theorem~\ref{thm:approx}. Without
uniqueness, the single-center grids of Section~\ref{sec:growth} can be
exponentially slow, and we show that unions of uniform cells restore a rate
that is polynomial for fixed bag size. The last subsection treats explicit
polynomial factors.

\subsection{Scaled data}\label{sec:exact-data}

Throughout Sections~\ref{sec:exact-data}--\ref{sec:exact-nonunique}, $F$ is a
quadratic with rational data,
\[
F(x)=\tfrac12x^{\T}Hx+b^{\T}x+c ,
\]
where $H$ is the symmetric Hessian, on a bounded mixed box
$X=\prod_iX_i$ with rational endpoints, after fixed coordinates have been
substituted and integer endpoints rounded inward. Write
$S=\argmin_XF$ and $\OPT=\min_XF$. Let $\Delta$ be a positive common denominator
of the monomial coefficients of $F$, namely $H_{ii}/2$, $H_{ij}$ ($i<j$),
$b_i$ and $c$, and of all endpoints $\ell_i,u_i$. Put
\begin{equation}\label{eq:exact-constants}
P=\Delta H,\qquad I_C^+=\{i\in I_C:H_{ii}>0\},\qquad
R=\Delta\prod_{i\in I_C^+}P_{ii},\qquad \Omega=\Delta R^2,\qquad \tau=\frac1{4nR}.
\end{equation}
The matrix $P$ is integral and symmetric, since $P_{ii}=2\Delta(H_{ii}/2)$ and
$P_{ij}=\Delta H_{ij}$; this choice of $\Delta$ absorbs the factor two of the
off-diagonal entries automatically. Each $P_{ii}$ with $i\in I_C^+$ is a
positive integer, and $P_{ii}\le\Delta L$ with $L=\max_iL_i$, so
$\log_2R\le\log_2\Delta+\abs{I_C^+}\log_2(\Delta L)$ is polynomial in $I$. The numbers
$R$, $\Omega$ and $\tau$ are computed from the original data; refined boxes never
enter them.

Two elementary facts are used repeatedly. If every coordinate of $x\in X$
lies in $\rho^{-1}\Z$ for a positive integer $\rho$, then
$\Delta\rho^2F(x)\in\Z$, because $\Delta$ clears every monomial coefficient. If
$a/W\ne a'/W'$ are rationals with positive denominators, then
$\abs{a/W-a'/W'}\ge1/(WW')$.

\begin{lemma}[Hadamard's inequality]\label{lem:hadamard}
If $A$ is a real symmetric positive definite $m\times m$ matrix, then
$\det A\le\prod_{i=1}^mA_{ii}$.
\end{lemma}

\begin{proof}
Let $\Lambda=\diag(A_{11},\dots,A_{mm})$, which is positive definite, and
$B=\Lambda^{-1/2}A\Lambda^{-1/2}$. Then $B$ is positive definite with unit
diagonal, so its eigenvalues $\lambda_k>0$ satisfy $\sum_k\lambda_k=m$, and
by the arithmetic--geometric mean inequality
$\det B=\prod_k\lambda_k\le1$. Hence
$\det A=\det B\prod_iA_{ii}\le\prod_iA_{ii}$.
\end{proof}

\subsection{Stationary polytopes and rational height}

\begin{definition}[Stationary polytope]\label{def:statpoly}
For $s\in S$ let $J_0(s)=\{i\in I_C:\ell_i<s_i<u_i\}$ and
\[
\mathcal P(s)=\bigl\{x\in\R^n:\ x_i=s_i\ (i\notin J_0(s)),\ \
\ell_i\le x_i\le u_i\ \text{and}\ \partial_iF(x)=0\ (i\in J_0(s))\bigr\}.
\]
\end{definition}

\begin{lemma}[Stationary polytopes]\label{lem:statpoly}
Let $s\in S$ and $J_0=J_0(s)$.
\begin{enumerate}[label=(\alph*)]
\item $H_{J_0J_0}\succeq0$.
\item $\mathcal P(s)$ is a nonempty polytope contained in $X$, and
$\mathcal P(s)\subseteq S$.
\item Let $v$ be a vertex of $\mathcal P(s)$ and
$T=\{i\in J_0:\ell_i<v_i<u_i\}$. Then $P_{TT}\succ0$, every coordinate of $v$
lies in $(\Delta\det P_{TT})^{-1}\Z$, and $1\le \Delta\det P_{TT}\le R$ (with
$\det P_{TT}=1$ if $T=\emptyset$).
\end{enumerate}
\end{lemma}

\begin{proof}
(a) For $w\in\R^n$ supported on $J_0$, the point $s+tw$ lies in $X$ for all
small $\abs t$, because the coordinates in $J_0$ are continuous and strictly
inside their bounds. Since $F(s+tw)=\OPT+t\nabla F(s)^{\T}w+\frac{t^2}2w^{\T}Hw$
has a minimum at $t=0$, we get $\nabla_{J_0}F(s)=0$ and $w^{\T}Hw\ge0$.

(b) By (a), $s\in\mathcal P(s)$. The set is defined by finitely many linear
equations and inequalities inside a bounded box, so it is a polytope. Its
points keep the integer coordinates of $s$ and satisfy all bounds, so
$\mathcal P(s)\subseteq X$. For $x\in\mathcal P(s)$, $d=x-s$ is supported on
$J_0$, and $H_{J_0J_0}d_{J_0}=\nabla_{J_0}F(x)-\nabla_{J_0}F(s)=0$. Hence
\[
F(x)-F(s)=\nabla F(s)^{\T}d+\tfrac12d^{\T}Hd
=\nabla_{J_0}F(s)^{\T}d_{J_0}+\tfrac12d_{J_0}^{\T}H_{J_0J_0}d_{J_0}=0 ,
\]
so $x\in S$.

(c) Suppose $w\in\R^T$ is nonzero and $H_{J_0T}w=0$. Extend $w$ by zero to
$\tilde w\in\R^n$. Then $v\pm t\tilde w\in\mathcal P(s)$ for small $t>0$: the
equations are preserved, the coordinates in $T$ have slack, and all other
coordinates are unchanged. This contradicts that $v$ is a vertex. So
$H_{J_0T}$ has trivial kernel. If $w\in\R^T\setminus\{0\}$ had
$w^{\T}H_{TT}w=\tilde w^{\T}H\tilde w=0$, positive semidefiniteness of
$H_{J_0J_0}$ would give $H_{J_0J_0}\tilde w_{J_0}=0$, that is
$H_{J_0T}w=0$; hence $H_{TT}\succ0$ and $P_{TT}=\Delta H_{TT}\succ0$. In
particular $H_{ii}>0$ for $i\in T$, so $T\subseteq I_C^+$.

The coordinates of $v$ outside $T$ are integers or endpoints, so they lie in
$\Delta^{-1}\Z$. The rows $i\in T$ of $\nabla F(v)=0$, multiplied by $\Delta^2$, read
\[
P_{TT}(\Delta v_T)=-\Delta(\Delta b_T)-\sum_{j\notin T}P_{Tj}(\Delta v_j),
\]
whose right-hand side is integral. Cramer's rule gives
$\Delta v_T\in(\det P_{TT})^{-1}\Z^T$, so every coordinate of $v$ lies in
$(\Delta\det P_{TT})^{-1}\Z$. Finally $\det P_{TT}$ is a positive integer, and by
Lemma~\ref{lem:hadamard} and $T\subseteq I_C^+$,
$\det P_{TT}\le\prod_{i\in T}P_{ii}\le\prod_{i\in I_C^+}P_{ii}$.
\end{proof}

\begin{corollary}[Rational height]\label{cor:height}
(a) Some global minimizer has all coordinates in $\rho^{-1}\Z$ for a positive
integer $\rho\le R$; every isolated global minimizer has this property.
(b) $\OPT=a/W$ in lowest terms with $1\le W\le \Omega$.
\end{corollary}

\begin{proof}
(a) Take any $s\in S$. The nonempty polytope $\mathcal P(s)$ has a vertex,
which is a global minimizer by Lemma~\ref{lem:statpoly}(b) and has the
stated height by Lemma~\ref{lem:statpoly}(c). If $s$ is isolated in $S$, then
the convex set $\mathcal P(s)\subseteq S$ contains no other point near $s$,
hence equals $\{s\}$, and $s$ is its vertex.
(b) Evaluate $F$ at the minimizer of (a): $\Delta\rho^2\OPT\in\Z$, so
$W\le \Delta\rho^2\le \Omega$.
\end{proof}

\begin{remark}[Other admissible height constants]\label{rem:heights}
Any bound on the minors of $P$ restricted to $I_C$ also works in
Lemma~\ref{lem:statpoly}(c) and Corollary~\ref{cor:height}, for example
$\abs{I_C}!\,(\max_{ij}\abs{P_{ij}})^{\abs{I_C}}$ by the Leibniz expansion, or
$\prod_{i\in I_C}\max\{1,\sum_{j\in I_C}\abs{P_{ij}}\}$ by Hadamard's
inequality for rows. Both also bound \emph{nonprincipal} minors, which arise
if a vertex of $\mathcal P(s)$ is described by a nonprincipal set of
stationarity rows. Part~(c) shows that a principal system $P_{TT}$ always
suffices, and Lemma~\ref{lem:hadamard} then gives the diagonal bound
in~\eqref{eq:exact-constants}, which is never larger than either
alternative and depends only on the coordinate curvatures. All these
constants have polynomial bit length; only the threshold values below
change.
\end{remark}

\subsection{Acceptance and recovery}

\begin{proposition}[Exact acceptance]\label{prop:accept}
Let $\beta\le\OPT$ and let $x\in X$ with $F(x)=a/W$ in lowest terms. If
$F(x)-\beta<1/(\Omega W)$, then $x\in S$.
\end{proposition}

\begin{proof}
If $F(x)>\OPT$, then $F(x)-\OPT\ge1/(\Omega W)$, since $\OPT$ has denominator at
most $\Omega$ by Corollary~\ref{cor:height}(b). With $\beta\le\OPT$ this
contradicts $F(x)-\beta<1/(\Omega W)$.
\end{proof}

The rule uses the denominator of the candidate itself. It does not assume
that feasible values lie on a lattice, which is false when continuous
coordinates vary.

\paragraph{Procedure REC$(y)$ (snapping recovery).}
Given $y\in X$: fix every integer coordinate at $y_i$; for every continuous
coordinate, fix it at $\ell_i$ if $y_i-\ell_i\le\tau$, else at $u_i$ if
$u_i-y_i\le\tau$, and otherwise call it free; let $J$ be the free set. Find
by exact linear programming a point $\hat x$ that agrees with the fixed
values, satisfies $\ell_J\le\hat x_J\le u_J$ and $\nabla_JF(\hat x)=0$, or
report failure. The data of this linear system are original coefficients,
endpoints and integer values of $y$, so it is solved in time polynomial in
$I$ \cite[Theorem~6.4.12]{GrotschelLovaszSchrijver1988}. If $H_{JJ}$ is
nonsingular, Gaussian elimination suffices.

\begin{lemma}[Snapping recovery]\label{lem:snap}
If $y\in X$ and $\dist(y,S)\le\tau/2$, then REC$(y)$ does not fail, and every
point it can return lies in $S$.
\end{lemma}

\begin{proof}
Let $s\in S$ be nearest to $y$; it exists since $S$ is compact. As
$\norm{y-s}\le\tau/2<1$, the integer coordinates of $y$ and $s$ agree. Two
distinct endpoints of a continuous coordinate differ by at least
$1/\Delta\ge4n\tau>2\tau$, so no coordinate is within $\tau$ of both endpoints. If
$s_i=\ell_i$ then $y_i-\ell_i\le\tau/2$ and coordinate $i$ is fixed at
$\ell_i$; similarly for $s_i=u_i$. Hence every continuous coordinate at a
bound in $s$ is fixed at its value in $s$, and $J\subseteq J_0=J_0(s)$. Let
$E=J_0\setminus J$ be the remaining fixed coordinates, each fixed at an
endpoint $e_i\in\{\ell_i,u_i\}$, and define on $\mathcal P(s)$ the linear
function
\[
\phi(x)=\sum_{i\in E}\abs{x_i-e_i}
=\sum_{i\in E,\,e_i=\ell_i}(x_i-\ell_i)+\sum_{i\in E,\,e_i=u_i}(u_i-x_i)\ge0 .
\]
For $i\in E$ we have $\abs{y_i-e_i}\le\tau$ and $\abs{s_i-y_i}\le\tau/2$, so
$\phi(s)\le\frac32\abs E\tau\le\frac3{8R}<\frac1R$. The minimum of $\phi$
over the polytope $\mathcal P(s)$ is attained at a vertex $v$. By
Lemma~\ref{lem:statpoly}(c) the coordinates of $v$ and all endpoints lie in
$\rho^{-1}\Z$ for some integer $\rho\le R$, so $\phi(v)\in\rho^{-1}\Z$.
Since $0\le\phi(v)\le\phi(s)<1/R\le1/\rho$, we get $\phi(v)=0$.

Thus $s'=v\in\mathcal P(s)\subseteq S$ satisfies $s'_i=e_i$ for $i\in E$,
agrees with $s$ (hence with the fixed values of REC) outside $J_0$, and has
$\nabla_JF(s')=0$ because $J\subseteq J_0$. So the system solved by REC is
feasible. If $\hat x$ is any solution, then $d=\hat x-s'$ is supported on
$J$ and $H_{JJ}d_J=\nabla_JF(\hat x)-\nabla_JF(s')=0$, so
$F(\hat x)-F(s')=\nabla_JF(s')^{\T}d_J+\frac12d_J^{\T}H_{JJ}d_J=0$.
Since $\hat x\in X$, it is a global minimizer. No nonsingularity of $H_{JJ}$
is used.
\end{proof}

\begin{theorem}[From certified gaps to exact minimizers]\label{thm:transfer}
Let $\beta\le\OPT$ be certified and $y\in X$.
\begin{enumerate}[label=(\alph*)]
\item If $F$ has set growth $F(x)-\OPT\ge g_S\dist(x,S)^2$ on $X$ with
$g_S>0$, and
\begin{equation}\label{eq:transfer-threshold}
F(y)-\beta\le\varepsilon_S:=\min\Bigl\{\frac{g_S\tau^2}{4},\
\frac1{2\Omega^2}\Bigr\}
=\min\Bigl\{\frac{g_S}{64n^2R^2},\ \frac1{2\Omega^2}\Bigr\},
\end{equation}
then REC$(y)$ returns a point $\hat x$ that passes the test of
Proposition~\ref{prop:accept}.
\item Without any growth assumption there is $\varepsilon_0>0$, depending on
the instance, such that the conclusion of (a) holds whenever
$F(y)-\beta\le\varepsilon_0$.
\end{enumerate}
\end{theorem}

\begin{proof}
(a) $g_S\dist(y,S)^2\le F(y)-\OPT\le F(y)-\beta\le g_S\tau^2/4$, so
Lemma~\ref{lem:snap} applies and $\hat x\in S$. Then
$F(\hat x)=\OPT=a/W$ with $W\le \Omega$, and
$F(\hat x)-\beta\le F(y)-\beta\le1/(2\Omega^2)<1/(\Omega W)$.
(b) The set $Y=\{x\in X:\dist(x,S)\ge\tau/2\}$ is compact and $F>\OPT$ on
it, so $\eta=\min_YF-\OPT>0$ (put $\eta=\infty$ if $Y=\emptyset$). With
$\varepsilon_0=\min\{\eta/2,1/(2\Omega^2)\}$, $F(y)-\beta\le\varepsilon_0$ gives
$\dist(y,S)<\tau/2$, and the argument of (a) applies.
\end{proof}

\begin{remark}[Set growth is automatic]\label{rem:setgrowth}
For every rational mixed box QP some $g_S>0$ exists. On an integer slice
whose continuous box contains a minimizer, Luo and Sturm's error bound for a
quadratic on a polytope \cite[Theorem~3.3]{LuoSturm2000}, applied to $F-\OPT$,
gives $F(x)-\OPT\ge c\,\dist(x,S\cap\text{slice})^2\ge c\,\dist(x,S)^2$; on
the finitely many other slices $F-\OPT$ is bounded below by a positive
constant while $\dist(x,S)$ is bounded by the diameter of $\bar X$. As for
point growth, no useful lower bound on $g_S$ follows. Theorem~\ref{thm:transfer}
shows that exactness costs only $\log_2(1/\varepsilon_S)
=\poly(I)+\log_2(1/g_S)$ bits of certified accuracy; the cost of reaching that
accuracy is a property of the approximation algorithm.
\end{remark}

\subsection{Exact output under quadratic growth}

\begin{lemma}[Uniqueness implies growth]\label{lem:unique-growth}
If $S=\{x^*\}$, then~\eqref{eq:growth} holds for some $g>0$.
\end{lemma}

\begin{proof}
Otherwise there are $x^k\in X\setminus\{x^*\}$ with
$(F(x^k)-\OPT)/\norm{x^k-x^*}^2\to0$. As $\bar X$ is bounded,
$F(x^k)\to\OPT$, and compactness of $X$ with uniqueness gives $x^k\to x^*$;
so the integer coordinates of $x^k$ equal those of $x^*$ for large $k$;
since $x^k\ne x^*$, they differ in some continuous coordinate. Put
$t_k=\norm{x^k-x^*}$ and
$u^k=(x^k-x^*)/t_k$; pass to a subsequence with $u^k\to u$, $\norm u=1$,
$u_{I_Z}=0$. Then
\[
\frac{F(x^k)-\OPT}{t_k^2}=\frac{\nabla F(x^*)^{\T}u^k}{t_k}
+\frac12(u^k)^{\T}Hu^k .
\]
Since $x^*$ minimizes $F$ over the convex continuous box of its integer slice
and $x^k$ lies in that box, the first term is nonnegative. The left side
tends to zero and the second term is bounded, so
$\nabla F(x^*)^{\T}u=0$ and $u^{\T}Hu\le0$. Each $u^k$ satisfies
$u^k_i\ge0$ when $x^*_i=\ell_i$ and $u^k_i\le0$ when $x^*_i=u_i$; so does
$u$, and hence $x^*+tu\in X$ for small $t>0$. Then
$F(x^*+tu)-\OPT=\frac{t^2}2u^{\T}Hu\le0$, contradicting uniqueness.
\end{proof}

The constant $g$ can be exponentially small in $I$.

\paragraph{Algorithm EX (exact output).}
For $q=1,2,4,8,\dots$: run CT$(2^{-q})$ and obtain a feasible $\hat x$ and a
certified $\beta$; run REC$(\hat x)$; if it returns $\tilde x$ with
$F(\tilde x)-\beta<1/(\Omega W)$, where $W$ is the reduced denominator of
$F(\tilde x)$, return $\tilde x$, $F(\tilde x)$ and the certificate of
CT. If $L=0$, CT is exact at its first stage and the test passes with
$\tilde x=\hat x$.

\begin{theorem}[Exact rational output]\label{thm:exact}
For every rational mixed box QP with a supplied tree decomposition, algorithm
EX stops, and its output is a global minimizer and $\OPT$, with a certificate
whose validity uses neither growth nor uniqueness. If the instance has
quadratic growth~\eqref{eq:growth} with condition number $\kappa$, EX uses
at most $f_1(p,\kappa)(I+1)^{C_1}$ bit operations, where $C_1$ is an absolute
constant; neither $g$ nor $\kappa$ is an input. This bound remains valid if
REC skips
every call whose matrix $H_{JJ}$ is singular, so under growth no linear
programming is needed.
\end{theorem}

\begin{proof}
\emph{Validity.} Each returned point passes Proposition~\ref{prop:accept}
with a certified $\beta$ (Theorem~\ref{thm:approx}), so it is a global
minimizer.

\emph{Termination.} CT stops on every instance
(Theorem~\ref{thm:approx}(a)). Once $2^{-q}\le\varepsilon_0$ of
Theorem~\ref{thm:transfer}(b), the test passes.

\emph{Bound under growth.} Point growth is set growth with $S=\{x^*\}$ and
$g_S=g$. By Theorem~\ref{thm:transfer}(a), EX stops at the first $q$ with
$2^{-q}\le\varepsilon_S$, so the largest $q$ used is less than
$2q^*$ with $q^*=\lceil\log_2(1/\varepsilon_S)\rceil$. Since $g\ge L/\kappa$
and $\log R$, $\log \Omega$, $\log n$, $\log(1/L)$ are at most polynomial in $I$,
$q^*\le\poly(I)+\log_2\kappa$. There are $O(\log q^*)$ calls, each within
$f(p,\kappa)(I+2q^*+1)^C$ bit operations by Theorem~\ref{thm:approx}, and
$(I+2q^*+1)^C\le(1+\log_2\kappa)^{C}\poly(I)$. REC, the evaluation of
$F(\tilde x)$ and the comparison cost $\poly(I+q^*)$. Absorbing all powers
of $1+\log\kappa$ into $f_1$ proves the bound. Finally, if
$\dist(\hat x,x^*)\le\tau/2$, then $J\subseteq J_0(x^*)$ as in
Lemma~\ref{lem:snap}, and $H_{J_0J_0}\succ0$: a nonzero $w$ with
$H_{J_0J_0}w=0$ would make $F$ constant on a short segment through $x^*$.
So $H_{JJ}$ is nonsingular at every call that the argument above needs, and
skipping singular calls does not change the bound.
\end{proof}

\begin{remark}[Coordinatewise reconstruction]\label{rem:cf}
Under uniqueness one may instead round each coordinate of $\hat x$ to the
unique rational of denominator at most $R$ within $1/(4R^2)$ (continued
fractions, \cite[Theorem~5.1.9]{GrotschelLovaszSchrijver1988}); this succeeds
once $2^{-q}\le g/(32R^4)$, by Corollary~\ref{cor:height}(a) and the
$1/R^2$ separation of such rationals. Snapping recovery needs only
$g/(64n^2R^2)$, does not need uniqueness, and returns a minimizer even when
$S$ is a continuum.
\end{remark}

\subsection{Explicit polynomial factors}\label{sec:polynomial}

Now let each factor be an explicit list of rational monomials of degree at
most $d\ge2$, with $d$ fixed (or encoded in unary).

\begin{lemma}[Interval curvature certificate]\label{lem:intcurv}
For each $i$ write $\partial_{ii}F=\sum_m\gamma_mx^{\alpha_m}$ and let
$[\underline\mu_m,\overline\mu_m]$ be the range of the monomial $x^{\alpha_m}$
on $\bar X$. Put $\hat L_i=\sum_m\max\{\gamma_m\overline\mu_m,
\gamma_m\underline\mu_m\}$ and $L_i=\max\{\hat L_i,0\}$. Then the $L_i$ are
upper coordinate curvatures, and they are computed in time polynomial in $I$.
\end{lemma}

\begin{proof}
The monomial $x^{\alpha}$ is a product of functions $x_k^{\alpha_k}$ of
distinct coordinates, so its range on the box $\bar X$ is the product of the
intervals $\{t^{\alpha_k}:t\in[\ell_k,u_k]\}$, each of which is computed
exactly from $\ell_k^{\alpha_k}$, $u_k^{\alpha_k}$ and, for even $\alpha_k$
with $\ell_k<0<u_k$, the value $0$. Hence
$\gamma_mx^{\alpha_m}\le\max\{\gamma_m\overline\mu_m,\gamma_m\underline\mu_m\}$
on $\bar X$, and $\partial_{ii}F\le\hat L_i\le L_i$. For every $x\in\bar X$,
$t\mapsto F(x+te_i)-\frac{L_i}2t^2$ has second derivative
$\partial_{ii}F-L_i\le0$ on its interval, so it is concave. The endpoints are
powers of degree at most $d$ of input rationals, so their bit length is
$O(dI)$, and there are at most as many terms as input monomials.
\end{proof}

The bound can overestimate when monomials cancel; the condition number in
the results below is computed with the accepted $L$. Any sharper bound with a
certificate checkable in polynomial time may be used instead.

\begin{proposition}[Lattice stopping rule]\label{prop:lattice}
Suppose every coordinate is integer, and let $Q_0$ be a positive common
denominator of the coefficients of the objective (after substituting fixed
coordinates). If $\hat x\in X$, $\beta\le\OPT$ and
$F(\hat x)-\beta<1/Q_0$, then $\hat x\in S$. It suffices that
$F(\hat x)-\beta\le2^{-q}$ with $q=1+\lceil\log_2Q_0\rceil=O(I)$.
\end{proposition}

\begin{proof}
At integer points every monomial value lies in $Q_0^{-1}\Z$, so
$F(X)\subseteq Q_0^{-1}\Z$. If $F(\hat x)>\OPT$ then
$F(\hat x)-\OPT\ge1/Q_0>F(\hat x)-\beta$, contradicting $\beta\le\OPT$.
\end{proof}

\begin{corollary}[Polynomial factors]\label{cor:poly}
For fixed $d$, under quadratic growth~\eqref{eq:growth} with the accepted
curvature bounds, CT$(2^{-q})$ returns a certified $2^{-q}$-approximation in
$f_d(p,\kappa)(I+q+1)^{C}$ bit operations with $C$ absolute. If all
coordinates are integer, CT$(2^{-(1+\lceil\log_2Q_0\rceil)})$ returns an exact
minimizer, certified by Proposition~\ref{prop:lattice}, in
$f_d(p,\kappa)(I+1)^C$ bit operations.
\end{corollary}

\begin{proof}
The interpolation, filtering, contraction, localization and grid-size
arguments of Sections~\ref{sec:grids} and~\ref{sec:growth} use only
Definition~\ref{def:curvature} and growth, not the quadratic form. In the bit
analysis of Theorem~\ref{thm:approx} only factor evaluation changes: a
monomial of degree at most $d$ at nodes with denominators dividing $A_j$ has
denominator dividing $\Delta_FA_j^d$, where $\Delta_F$ clears the coefficients, and
if the nodes have bit length at most $\lambda$ its value has bit length
$O(I+d\lambda)$; corrections and
messages then have a common denominator dividing $8\Delta_FA_j^{d}$ (with $\Delta_F$
also clearing $L$). Evaluation of an explicit factor costs a polynomial in
its length and these bit lengths, for fixed $d$. The second statement
follows from Proposition~\ref{prop:lattice} with $q=O(I)$.
\end{proof}

\begin{example}[Limits of the polynomial extension]\label{ex:polylimits}
(a) $F(x)=x^3-6x$ on $[1,2]$ has $\partial_{xx}F=6x\le12=\hat L$ and
$F(x)-\OPT=(x-\sqrt2)^2(x+2\sqrt2)\ge3(x-\sqrt2)^2$, so it has growth $g=3$
and $\kappa=4$; but its minimizer $\sqrt2$ and value $-4\sqrt2$ are
irrational, so no rational exact output exists for continuous polynomial
coordinates. (b) $F(x)=x^4$ on $[-1,1]$ has a unique minimizer and no
$g>0$, so Lemma~\ref{lem:unique-growth} fails for polynomials.
(c) For $k\ge1$, $F(x,y)=x^2-\frac12x^{2^k}+y^2$ on $[0,1]^2$ has
$\partial_{xx}F\le2$, $\partial_{yy}F=2$ and growth $g=\frac12$, since
$x^{2^k}\le x^2$; its input has $O(k)$ bits if exponents are written in
binary, but $F(\frac12,0)$ has denominator $2^{2^k+1}$. Binary-encoded
exponents therefore need a separate evaluation analysis.
\end{example}

For continuous polynomial objectives an exact answer must therefore be
implicit. Appendix~\ref{app:boundary} gives one form: under point growth and
strict complementarity at the active bounds, sign tests on the certified
retained box identify the active face, and a uniform convexity test on the
remaining coordinates yields a strongly convex polynomial subproblem whose
unique minimizer is the global one. The cost carries a term
$B_\gamma=\lceil\log_2\max\{2,C_2/\gamma\}\rceil$ in the smallest active
derivative $\gamma$; Propositions~\ref{prop:margin} and~\ref{prop:weakcompl} show
that this term is intrinsic to certificates of this form and that weak
complementarity defeats them.
